% !PRINTED 420
\weberpage{420}
%<<CONT
$X_b$ theilbar sein, und die $\beta^a$ und folglich auch die $\beta$ selbst,
die in (6) vorkommen, \gesperrt{umfassen zusammen alle primitiven $b^{\mathrm{ten}}$
Einheitswurzeln.} Ebenso wird nun bewiesen, dass $\alpha$ in (6) alle
primitiven $a^{\mathrm{ten}}$ Einheitswurzeln durchläuft, und dass also
$\rho$ alle primitiven $n^{\mathrm{ten}}$ Einheitswurzeln durchläuft.

Die Wurzeln von $\varphi(x)$ stimmen also völlig überein mit den sämmtlichen
Wurzeln von $X_n$, und folglich muss, wenn wir zur Bestimmung eines
Zahlenfactors etwa noch festsetzen, dass die höchste Potenz von $x$ in
$\varphi(x)$ den Co\"efficienten $1$ haben soll, $\varphi(x)$ mit $X_n$
\gesperrt{identisch sein}.

\gesperrt{Es ist also unter den gemachten Voraussetzungen $X_n$ nicht in Factoren
niedrigeren Grades mit rationalen Co\"efficienten zerlegbar.}

Mit denselben Hülfsmitteln können wir den folgenden Satz von
\gesperrt{Kronecker} beweisen, der die Irreducibilität der Functionen $X_n$
in einem weiteren Sinne ausspricht\footnote{Mém. sur les facteurs
irréductibles de l'expression $(x^n - 1)$. Liouville's Journal, Vol. XIV.}.

\gesperrt{Die Function $X_n$ kann nicht in Factoren zerlegt werden, die in
Bezug auf $x$ ganz und rational sind, und in den Co\"efficienten nur
rationale Zahlen und Einheitswurzeln enthalten, deren Grad relativ prim zu
$n$ ist.}

Angenommen, es sei $a$ relativ prim zu $n$, und $r$ eine primitive
$n^{\mathrm{te}}$, $\alpha$ eine primitive $a^{\mathrm{te}}$ Einheitswurzel,
und es sei $\varphi(x,\alpha)$ ein Theiler von $X_n$, also wenigstens für
ein $r$

\weq{7}{\varphi(r,\alpha) = 0.}

Es sei nun $\rho$ eine primitive $na^{\mathrm{te}}$ Einheitswurzel, also
eine Wurzel von $X_{na} = 0$; dann ist sowohl $r$ als $\alpha$ eine Potenz
von $\rho$ (\S.\,131), etwa
\[
r = \rho^{\xi},\quad \alpha = \rho^{\eta},
\]
worin $\xi$ relativ prim zu $n$, aber theilbar durch $a$, $\eta$ relativ
prim zu $a$, aber theilbar durch $n$ ist. Es ist also

\weq{8}{\varphi(\rho^{\xi}, \rho^{\eta}) = 0.}

Nun ist, wie wir soeben bewiesen haben, $X_{na}$ irreducibel, und folglich
muss die Function $\varphi(x^{\xi}, x^{\eta})$, die mit $X_{na}$ einen
Theiler gemein hat, durch diese Function theilbar sein, d. h. es ist für
jedes ganzzahlige $h$, was zu $na$ relativ prim ist:

\weq{9}{\varphi(\rho^{h\xi}, \rho^{h\eta}) = 0.}
